\section{Related Work}
\label{sec:related-work}

% PROPOSED rebuild for IST/JSS — draft artifact (does not replace manuscript until integrated).
% Bibliography: related_work_rebuild/related_work.bib (+ grey_literature.bib).
% All scholarly citations below are Crossref-verified unless noted.

Empirical software engineering increasingly discovers GitHub repositories by searching for agent-instruction files (\texttt{AGENTS.md}, editor rule sets, Copilot instruction paths, prompt directories, and related filenames).
A satisfied path predicate is a retrieval event, not a membership proof for the analytic population named in a research question.
We call a repository \emph{contaminated} when it satisfies the discovery search yet falls outside that stated population.
This section situates prior work along six clusters and ends at the gap our protocol addresses.
We deliberately distinguish three contamination senses that the literature often conflates: (i)~sample--target mismatch (our construct), (ii)~train/test or duplication leakage in learned evaluations, and (iii)~label-process failure inside an otherwise well-scoped dataset.

\subsection{MSR methodology, discovery infrastructure, and GitHub affordances}
\label{subsec:rw-msr}

Mining software repositories was articulated as a research programme with explicit methodological risks~\cite{hassan2008road}.
Forge mining warnings predate GitHub: SourceForge status and metadata systematically mislead naive samples~\cite{howison2004perils}.
On GitHub, \textsc{GHTorrent} and its lean on-demand variant made event streams and repository metadata routinely available~\cite{gousios2012ghtorrent,gousios2014lean}; PyDriller lowered the cost of commit-level extraction~\cite{spadini2018pydriller}; World of Code and Software Heritage expanded discovery beyond a single forge's search API~\cite{ma2019woc,pietri2019swh,dicossmo2018swh}.
A systematic mapping of GitHub-based software engineering research catalogues topics and methods but treats discovery largely as an access problem~\cite{cosentino2017mapping}.

These lines largely solve \emph{coverage and tooling}.
They do not, by themselves, establish that repositories retrieved by instruction-file predicates belong in a study-specific analytic population, nor do they report how contamination estimates move when membership labels are aggregated under alternative consensus protocols.

\subsection{Sampling methodology, GitHub perils, and curation cousins}
\label{subsec:rw-sampling}

Kalliamvakou et al.\ document that GitHub entities are often not software projects, that activity proxies are fragile, and that mining pipelines inherit those mismatches~\cite{kalliamvakou2014perils,kalliamvakou2015indepth}.
Baltes and Ralph critically review sampling practice in software engineering and argue for explicit frames, populations, and generalization claims~\cite{baltes2022sampling}.
Dabic et al.\ study how MSR papers sample GitHub projects and how sampling choices shape corpora~\cite{dabic2021sampling}.
Stars, social signals, and related metadata are themselves unstable popularity and quality proxies~\cite{borges2016popularity,tsay2014influence}.

A related but distinct literature \emph{curates} GitHub for engineered software projects.
RepoReapers classifies engineered versus non-engineered repositories~\cite{munaiah2017curating}; PHANTOM revisits curation with time-series clustering~\cite{phantom2020}; Golzadeh et al.\ build ground truth for another membership class---bots~\cite{golzadeh2021bots}.
These works supply transferable lessons (worksheeted labels, multi-signal features, held-out validation), yet their constructs are not ours: engineered status (or bot status) is neither necessary nor sufficient for membership in an instruction-artifact analytic population (for example, an engineered AI-builder SDK can be on-target for a tooling study and off-target for a conventional-application study).

\textbf{Gap relative to this article.}
Prior sampling and curation work motivates auditing sample--target fit but does not quantify instruction-artifact contamination by predicate family or under complementary consensus protocols.

\subsection{Construct validity, threats to validity, and reporting guidelines}
\label{subsec:rw-validity}

Classical experimental and case-study guidance treats operational definitions, constructs, and reporting as first-class validity objects~\cite{wohlin2012exp,runeson2009case,jedlitschka2005reporting}.
Siegmund et al.\ show that the community itself disagrees about the relative priority of internal and external validity~\cite{siegmund2015views}, which makes undeclared analytic targets especially dangerous.
Ampatzoglou et al.\ taxonomize and mitigate threats to validity in secondary studies~\cite{ampatzoglou2019threats}; Verdecchia et al.\ critically reflect on how threats are discussed in software engineering research more broadly~\cite{verdecchia2023threats}.
Systematic mapping guidelines and the ACM SIGSOFT empirical standards announcement further pressure authors to report design choices~\cite{petersen2015mapping,ralph2021standards}.

\textbf{Gap relative to this article.}
This validity literature explains \emph{why} inclusion rules and operational definitions matter; it does not provide an operational discover--filter--annotate--inspect worksheet for AI-instruction discovery frames, nor empirical evidence that consensus-protocol choice shifts contamination rates on the same repositories.

\subsection{Dataset contamination, leakage, duplication, and label validation}
\label{subsec:rw-contamination}

``Contamination'' is overloaded.
In data mining, Kaufman et al.\ formalize train/test leakage~\cite{kaufman2012leakage}; Moreno-Torres et al.\ unify dataset-shift phenomena~\cite{moreno2012shift}; Kapoor and Narayanan document leakage-driven reproducibility failures across machine-learning science~\cite{kapoor2023leakage}.
On GitHub code corpora, D{\'e}j{\`a}Vu maps massive duplication~\cite{lopes2017dejavu}, and Allamanis shows that duplication inflates learned evaluations~\cite{allamanis2019duplication}.
Cross-project defect prediction classic results show that domain mismatch collapses performance even when retrieval is successful~\cite{zimmermann2009cross}.

A parallel dataset-validation line studies corrupted or process-fragile labels: NASA defect data quality~\cite{shepperd2013nasa}, researcher/process bias in defect prediction~\cite{shepperd2014bias}, tangled changes that break change-based constructs~\cite{herzig2013tangled}, SZZ and feature-collection failures under manual validation~\cite{herbold2022szz}, and sensitivity of conclusions to model-validation technique~\cite{tantithamthavorn2017}.
Benchmark construction work likewise insists on explicit scope~\cite{dambros2012benchmark}; SWE-bench illustrates large-scale GitHub-issue task construction with filtering choices that shape the evaluation population~\cite{jimenez2024swebench}.

\textbf{Gap relative to this article.}
These papers solve adjacent problems---evaluation leakage, duplication, and label-process failure---but not \emph{path-search population mismatch} for instruction-artifact frames.
We reuse their lesson that unvalidated denominators bias downstream claims, while keeping the estimand distinct.

\subsection{Reproducibility, replication packages, and research artifacts}
\label{subsec:rw-repro}

Repository-based studies are hard to replay when retrieval, filters, and labels are underspecified~\cite{gonzalez2011repro,gonzalez2023revisit}.
Artifact availability remains uneven~\cite{heumuller2020artifacts}; artifact evaluation practice has matured over a decade~\cite{winter2022artifacts}; recent surveys update status and trends~\cite{liu2024artifacts}; mapping work charts replication practice~\cite{silva2014replication}.

\textbf{Gap relative to this article.}
Reproducibility norms justify releasing worksheets, frozen labels, and alternative-protocol replay scripts; they do not substitute for the contamination audit itself.

\subsection{Human annotation, consensus, and LLM-assisted labeling}
\label{subsec:rw-annotation}

Chance-corrected agreement coefficients and interpretation benchmarks remain the reporting baseline for categorical coding~\cite{cohen1960kappa,landis1977kappa}.
LLM-assisted annotation is now common: Gilardi et al.\ report strong ChatGPT performance against crowd workers on political text~\cite{gilardi2023chatgpt}; Zheng et al.\ evaluate LLM-as-a-judge reliability~\cite{zheng2023judge}; Ahmed et al.\ ask whether LLMs can replace manual annotation of software engineering artifacts and give a more cautious, SE-specific answer~\cite{ahmed2025llms}.
Human--AI interaction guidelines caution against unsupervised automation of consequential judgments~\cite{amershi2019hai}.

Evidence is conflicting across domains: optimistic general-NLP results do not automatically transfer to repository membership boundaries, where metadata sparsity and product-role ambiguity concentrate disagreement.
Our protocol therefore treats consensus rules as part of the estimand and keeps human coding primary for membership labels.

\textbf{Gap relative to this article.}
Annotation research improves labeling technology; it does not measure how alternative consensus protocols change instruction-frame contamination rates, nor how disagreement clusters at membership boundaries versus product-role classes.

\subsection{AI instruction artifacts, promptware, and coding assistants}
\label{subsec:rw-instruction}

Promptware and LLM-for-SE surveys establish that prompts, agent memory files, and tool interfaces are becoming software engineering objects~\cite{chen2026promptware,hou2024llm4se,fan2023llmse}.
Empirical studies of Copilot and related assistants characterize productivity and interaction~\cite{ziegler2024copilot,barke2023grounded,vaithilingam2022expectation}; earlier bot work anticipated automation in developer workflows~\cite{storey2016bots}.
Primary documentation specifies the concrete filenames and configuration surfaces that mining studies now search for~\cite{agentsmd2025,brown2024copilot,claude2024docs,cursor2024docs,anthropic2024mcp}.

\textbf{Gap relative to this article.}
This cluster documents the phenomenon and, increasingly, its engineering.
It rarely audits whether repositories retrieved via those filenames match the analytic population a study claims, and it does not provide family-level contamination reporting or consensus-protocol sensitivity analyses.

\subsection{Synthesis: where this article begins}
\label{subsec:rw-synthesis}

What is already solved is awareness that GitHub samples are noisy, that sampling frames must be declared, that engineered-project filters help for some questions, that datasets need validation, and that artifacts should be replayable.
What is only partially solved is operational sample--target auditing for new discovery predicates, especially when membership is coder-mediated.

What remains unaddressed---and where our contribution begins---is an integrated audit for \emph{AI-instruction-artifact discovery frames} that (1)~defines contamination as target-conditional sample--target mismatch, (2)~reports complementary consensus estimates on the same repositories, (3)~decomposes misalignment by predicate family (and metadata sparsity), (4)~separates membership-boundary disagreement from product-role classification in a tiered inspection, and (5)~releases a worksheet protocol with frozen labels and replay scripts as reporting guidance.
No verified peer-reviewed study in our search closes that intersection; claiming otherwise would overstate novelty of the perils insight and understate novelty of the operational audit.
